\chapter{Kết luận và nhận xét tổng thể}

\section{Bảng tổng hợp toàn bài}

\begin{table}[H]
\centering
\caption{Tổng hợp kết quả tốt nhất trên ba bộ dữ liệu}
\label{tab:overall-project-summary}
\begin{tabularx}{\textwidth}{l l l c >{\raggedright\arraybackslash}X}
\toprule
Bộ dữ liệu & ANN tốt nhất & Chỉ số chính & Kết quả & Insight chính \\
\midrule
MNIST & Single-branch ANN & Test accuracy & 99.14\% & Chuẩn hóa hình học và HOG giải quyết phần lớn khó khăn; augmentation cục bộ và two-branch không tạo thêm lợi ích rõ rệt. \\
FashionMNIST & Single-branch ensemble & Test accuracy & 91.68\% & HOG và tiền xử lý vẫn hữu ích, nhưng bước cải thiện tốt nhất là ensemble để giảm phương sai dự đoán giữa các mô hình đơn. \\
ChestMNIST & Feature-fusion lite ANN & Macro AUC / ACC & 0.7424 / 0.7301 & Weighted loss, cross-validation ensemble và feature fusion là ba thành phần quyết định để ANN cạnh tranh được trên dữ liệu y sinh đa nhãn. \\
\bottomrule
\end{tabularx}
\end{table}


\section{Đóng góp chính của bài}

Trong phạm vi yêu cầu của Lab 03, báo cáo này có bốn đóng góp chính. Thứ nhất, báo cáo đã thiết kế và triển khai đầy đủ các pipeline ANN cho ba nhóm dữ liệu khác nhau: MNIST Handwritten, FashionMNIST và ChestMNIST. Thứ hai, mỗi bài toán đều không dừng ở một mô hình duy nhất mà được phát triển theo chuỗi baseline $\rightarrow$ cải tiến $\rightarrow$ đánh giá lại, nhờ đó có thể chỉ ra khá rõ thành phần nào thực sự tạo ra cải thiện. Thứ ba, báo cáo đã phân tích kết quả ở cả mức chỉ số tổng hợp lẫn mức lỗi điển hình, từ đó sinh ra các biến thể kế tiếp trên cơ sở thực nghiệm thay vì thử ngẫu nhiên. Thứ tư, báo cáo đã rút ra được các hướng cải tiến cụ thể cho từng bộ dữ liệu, thay vì chỉ kết luận chung rằng ``mô hình tốt hơn''.

\section{Kết luận thực nghiệm}

Ba nhóm thí nghiệm trên MNIST, FashionMNIST và ChestMNIST cho thấy cùng một xu hướng khá nhất quán: ANN thuần vẫn có thể đạt kết quả tốt nếu đầu vào được tổ chức hợp lý, nhưng hiệu quả của ANN phụ thuộc rất mạnh vào tiền xử lý, đặc trưng hóa và chiến lược tối ưu. Ở cả ba bộ dữ liệu, baseline đơn giản nhất luôn đóng vai trò mốc so sánh cần thiết để trả lời đúng câu hỏi đầu tiên của đề bài: ANN thuần trên raw pixel đi được tới đâu trước khi thêm các kỹ thuật khác.

Với \textbf{MNIST}, kết luận thực nghiệm rõ nhất là: cấu hình \textit{single-branch ANN} là lựa chọn tốt nhất với test accuracy $99.14\%$. Bước cải thiện quyết định nằm ở \textit{deslant + recenter + HOG}; trong khi đó, augmentation cục bộ và two-branch không tạo thêm lợi ích đáng kể. Vì vậy, với một bộ dữ liệu sạch như MNIST, ANN thuần hoàn toàn đủ mạnh nếu được cung cấp đúng biểu diễn đầu vào.

Với \textbf{FashionMNIST}, kết luận quan trọng nhất là: khó khăn chính không còn nằm ở việc thiếu đặc trưng toàn cục, mà nằm ở độ chồng lấn hình dáng giữa các lớp áo. Trong bối cảnh đó, cấu hình tốt nhất lại là \textit{ensemble của ba mô hình single-branch ANN} với $91.68\%$, chứ không phải two-branch hay region-focused feature fusion. Nói cách khác, trên FashionMNIST, bước cải thiện hiệu quả nhất là giảm phương sai dự đoán bằng ensemble.

Với \textbf{ChestMNIST}, cách chấm phải bám theo benchmark chuẩn là AUC và ACC. Kết quả tốt nhất hiện tại là \textit{feature-fusion lite ANN} với macro AUC $0.7424$ và ACC $0.7301$. Điều này cho thấy trên dữ liệu y sinh đa nhãn, ANN chỉ trở nên cạnh tranh khi đi cùng weighted loss, cross-validation ensemble và feature engineering phù hợp với texture cũng như cấu trúc mức xám của ảnh X-quang.

Nhìn ở mức tổng quát, bài học lớn nhất của toàn bộ báo cáo là: ANN không hề vô dụng trên dữ liệu ảnh, nhưng muốn ANN hiệu quả thì không thể chỉ tăng độ sâu hoặc độ rộng của MLP một cách cơ học. Điều quyết định hơn nằm ở việc tổ chức đầu vào đúng, chọn loss đúng và đánh giá đúng theo bản chất của từng bộ dữ liệu.

\section{Hạn chế cốt lõi}

Tuy vậy, báo cáo hiện tại vẫn có một giới hạn quan trọng. Tất cả các mô hình trong Lab 03 đều thuộc họ ANN thuần theo kiểu \textit{fully-connected}, nên chúng xem ảnh chủ yếu như một vector đặc trưng đã được trải phẳng. Cách làm này giúp việc phân tích và thiết kế feature engineering trở nên rõ ràng, nhưng đồng thời bỏ qua một giả định rất quan trọng của dữ liệu ảnh: các pixel gần nhau trong không gian thường có quan hệ chặt chẽ với nhau, còn các mẫu hình cục bộ như cạnh, góc, texture hay cấu trúc biên thường lặp lại ở nhiều vị trí khác nhau.

Đây chính là lý do sâu xa khiến ANN thuần trong báo cáo phải dựa nhiều hơn vào tiền xử lý, handcrafted feature và ensemble để bù lại những gì bản thân kiến trúc chưa nắm bắt được. Cũng vì giới hạn này mà khoảng cách với các benchmark mạnh hơn vẫn còn tồn tại. Trên FashionMNIST, benchmark công khai vẫn có những phương pháp vượt mốc hiện tại của báo cáo~\cite{fashionmnist_repo}. Trên ChestMNIST, dù cấu hình tốt nhất của báo cáo đã tiệm cận baseline AutoKeras, kết quả vẫn còn cách một khoảng đáng kể so với các baseline CNN mạnh hơn trong benchmark chính thức của MedMNIST~\cite{medmnist_site}.

\section{Đề xuất cải tiến theo từng bộ dữ liệu}

Với \textbf{MNIST}, hướng cải tiến hợp lý nhất không phải là thêm nhiều kỹ thuật mới, mà là giữ cấu hình single-branch ANN như một mốc đủ tốt rồi dùng nó làm baseline sạch cho các bài toán khó hơn. Kết quả của chương này đã cho thấy rất rõ rằng khi dữ liệu đủ sạch và đã được chuẩn hóa hình học, augmentation cục bộ và two-branch không còn tạo thêm giá trị đáng kể. Vì vậy, nếu tiếp tục mở rộng trên MNIST thì hướng phù hợp hơn là kiểm tra tính tối giản của mô hình, chẳng hạn giảm số tầng hoặc giảm số chiều ẩn để tìm ra cấu hình gọn hơn nhưng vẫn giữ vùng $99\%$.

Với \textbf{FashionMNIST}, điểm nghẽn hiện tại nằm ở nhóm lớp có silhouette rất gần nhau như \textit{T-shirt/top}, \textit{Shirt}, \textit{Pullover} và \textit{Coat}. Vì ensemble đang là bước cải thiện hiệu quả nhất, hướng phát triển tự nhiên là giữ single-branch ensemble làm trục chính, rồi thử những đặc trưng cục bộ mạnh hơn đúng ở vùng thân trên, thay vì tiếp tục mở rộng MLP một cách thuần túy. Những hướng đáng kiểm chứng hơn là spatial pyramid HOG, contour-based descriptor hoặc patch-based MLP; mục tiêu của chúng là giữ được thông tin \textit{vị trí} của cổ áo, vai và tay áo, vốn là phần mà vector pixel phẳng hoặc HOG toàn ảnh còn mô tả chưa đủ tách bạch~\cite{fashionmnist_repo}.

Với \textbf{ChestMNIST}, kết quả hiện tại đã cho thấy ANN có thể đi khá xa nếu đầu vào và loss được thiết kế đúng, nhưng khoảng cách với benchmark CNN vẫn còn. Vì vậy, bước tiếp theo hợp lý nhất là mở rộng từ feature-fusion lite ANN sang các mô hình vẫn giữ tinh thần MLP nhưng biết khai thác cấu trúc không gian tốt hơn, chẳng hạn patch-based MLP hoặc các khối attention nhẹ trên vùng ảnh. So với MNIST và FashionMNIST, đây là bài toán mà hạn chế ``vector hóa ảnh'' của ANN thuần lộ ra rõ nhất, nên nếu muốn tiếp tục tăng AUC thì việc đưa quan hệ không gian cục bộ trở lại mô hình sẽ là hướng phát triển tự nhiên hơn là chỉ tăng thêm chiều rộng của tầng ẩn~\cite{medmnist_site}.
